\documentclass[12pt]{article}
\usepackage{amsfonts,amsthm,amsmath,amssymb,mathtools}
\usepackage{mathrsfs}
\usepackage{enumitem}
\usepackage{tikz-cd}
\usepackage{geometry}
\usepackage{graphicx}
\IfFileExists{mlmodern.sty}{
    \usepackage{mlmodern}
}{}

\geometry{letterpaper, portrait, margin=1in}

\setcounter{section}{0}

\renewcommand{\thesection}{\S\arabic{section}}
\renewcommand{\thesubsection}{\arabic{section}}
\DeclareMathOperator{\im}{im}
\DeclareMathOperator{\id}{id}

\newtheorem{thm}{Theorem}
\newenvironment{theorem}[1]{
	\renewcommand{\thethm}{#1}
	\begin{thm}
		}{\end{thm}}

\newtheorem{lma}{Lemma}
\newenvironment{lemma}[1]{
	\renewcommand{\thelma}{#1}
	\begin{lma}
		}{\end{lma}}

\theoremstyle{definition}
\newtheorem{ex}{Problem}
\newenvironment{exercise}[1]{
	\renewcommand{\theex}{#1}
	\begin{ex}
		}{\end{ex}}

\title{Functional Analysis --- Homework 1}
\author{Connor Mooneyhan}
\date{}

\begin{document}

\maketitle

\begin{ex}
	For $p \in [1, \infty)$, give an example of a sequence which is in $\ell^q$ for all $q > p$, but not in $\ell^p$.
\end{ex}
\begin{proof}
	Consider the sequence $(x_n)$ defined by \begin{equation*}
		x_n = \left( \frac{1}{n} \right)^{\frac{1}{p}}.
	\end{equation*}
	Then $||(x_n)||_p^p = \sum_1^\infty 1/n$, which is the harmonic series, so $(x_n) \notin \ell^p$, but for any $q > p$ we have $(1/n)^{q/p} < 1/n$ when $n \neq 1$.
	Thus $||(x_n)||_q^q$ converges and $(x_n) \in \ell^q$.
\end{proof}

\begin{ex}
	If a subspace $S$ of a normed space $V$ contains a nonempty open set, then $S = V$.
\end{ex}
\begin{proof}
	Let $U$ be a nonempty open subset of $S$, let $p \in U$, and let $\delta > 0$ such that $B_\delta(p) \subset U$.
	Now consider an arbitrary $x \in V$ and $\alpha > 0$.
	The idea is to ``nudge p in the direction of $x$'' while remaining in $B_\delta(p)$ by finding $\alpha$ such that $\alpha(x-p) + p$ is close enough to $p$.
	Indeed, \begin{equation*}
		||[\alpha(x-p) + p] - p|| = ||\alpha(x-p)|| = \alpha||x-p||,
	\end{equation*}
	so by choosing $\alpha < \delta/||x-p||$, we get $\alpha(x-p) + p \in B_\delta(p) \subset S$.
	We can now use the closure of $S$ under addition and scalar multiplication to get that \begin{equation*}
		[\alpha(x-p) + p] + [\alpha p - p] = \alpha x \in S,
	\end{equation*}
	and thus multiplying by $1/\alpha$ we get $x \in S$.
	Since $x$ was arbitrary, we have shown that $V \subset S$, and thus $V = S$.
\end{proof}

\begin{ex}
	Show that in any normed space $V$, the operations \begin{equation*}
		+ : V \times V \to V \hspace{2em} \cdot : \mathbb{R} \times V \to V
	\end{equation*}
	of addition and scalar multiplication are continuous.
\end{ex}
\begin{proof}
	First we do scalar multiplication.
	Let $(x_n) \to x$ in $V$ and $(c_n) \to c$ in $\mathbb{R}$.
	Then if

	Now let $(x_n) \to x$ and $(y_n) \to y$ in $V$.
	Then given $\varepsilon > 0$, there exists $N \in \mathbb{N}$ (take such an integer for $(x_n)$ and one for $(y_n)$ and then $N$ is the max of these two) such that when $n > N$, \begin{align*}
		\frac{\varepsilon}{2} & > ||x - x_n|| \geq ||x|| - ||x_n||  \\
		\frac{\varepsilon}{2} & > ||y - y_n|| \geq ||y|| - ||y_n||.
	\end{align*}
	\textcolor{red}{UNFINISHED}
\end{proof}

\begin{ex}
	Let $V$ be a normed space.
	Show that a linear map $f : V \to \mathbb{R}$ is continuous if and only if $\ker(f)$ is closed.
\end{ex}
\begin{proof}
	If $f$ is continuous, then since $\ker(f)$ is the preimage of the closed set $\lbrace 0 \rbrace$, $\ker(f)$ is closed.
	If $\ker(f)$ is closed, let $(x_n)$ be a sequence converging to $x$ in $V$.

	\dots Thus $(f(x_n)) \to f(x)$, so $f$ is continuous as desired.
	\textcolor{red}{UNFINISHED}

\end{proof}

\begin{ex}
	Let $\rho : V \to [0, +\infty)$ be a seminorm.
	Prove that if there exists a positive number $\delta > 0$ such that the preimage $\rho^{-1}[0, \delta]$ contains an open ball, then $\rho$ is continuous.
\end{ex}
\begin{proof}
	\textcolor{red}{UNFINISHED}
\end{proof}

\begin{ex}
	Let $V$ be a normed space.
	Show that $V$ is Banach if and only if every absolutely convergent series is convergent.
\end{ex}
\begin{proof}
	First, assume $V$ is Banach and $\sum ||x_i|| < \infty$.
	Then $(||x_i||)$ is Cauchy in $\mathbb{R}$, so given $\varepsilon > 0$, we can choose $N > 0$ such that when $m > n > N$, \begin{equation*}
		\sum_1^m||x_i|| - \sum_1^n||x_i|| = \sum_{n+1}^m||x_i|| < \varepsilon.
	\end{equation*}
	Then we also have that \begin{equation*}
		\left|\left|\sum_1^m x_i - \sum_1^n x_i\right|\right| = \left|\left|\sum_{n+1}^m x_i\right|\right| \leq \sum_{n+1}^m ||x_i|| < \varepsilon,
	\end{equation*}
	so the partial sums of $\sum_1^\infty x_i$ are also Cauchy and thus converge since $V$ is Banach.

	Now if every absolutely convergent series is convergent, then let $(x_i)$ be a Cauchy sequence in $V$.
	Define $y_1 \coloneq x_1$, and then recursively define $y_n \coloneq x_n - \sum_{i=1}^{n-1}y_i$ for all $n > 1$, so that $x_n = \sum_1^n y_i$ for all $n$.
	We then just need to show that $\sum y_i$ is absolutely convergent, and we will have shown that $x_n$ converges and thus $V$ is Banach.

	\textcolor{red}{UNFINISHED}
\end{proof}

\begin{ex}
	In this problem, you will show that for any number $p > 1$, the function \begin{equation*}
		||x||_p = ||(x_1, \dots, x_n)||_p = (|x_1|^p + \cdots + |x_n|^p)^{1/p}
	\end{equation*}
	is a norm on $\mathbb{R}^n$.
	The extension to $\ell^p$ follows easily from that.
	It will be helpful to define a number $q$, called the conjugate exponent to $p$, by the equation: $\frac{1}{p} + \frac{1}{q} = 1$.
	\begin{enumerate}[label=(\alph*)]
		\item (Young's inequality) Use the fact that $\ln(x)$ is a concave function to prove that if $a, b$ are positive numbers, then \begin{equation*}
			      ab \leq \frac{a^p}{p} + \frac{b^q}{q}.
		      \end{equation*}
		\item (Hölder's inequality) Given any $x, y \in \mathbb{R}^n$, use Young's inequality to prove that \begin{equation*}
			      |\langle x, y \rangle| \leq ||x||_p||y||_q
		      \end{equation*}
		      where $\langle x, y \rangle$ denotes the standard dot product in $\mathbb{R}^n$.
		\item (Minkowski's inequality) Given any $x, y \in \mathbb{R}^n$, use Hölder's inequality to prove \begin{equation*}
			      ||x+y||_p \leq ||x||_p + ||y||_p,
		      \end{equation*}
		      which is the triangle inequality for $||\cdot||_p$.
	\end{enumerate}
\end{ex}
\begin{proof}\
	\begin{enumerate}[label=(\alph*)]
		\item Since $\ln(x)$ is concave and the tangent line at $x_0$ is \begin{align*}
			      y & = \frac{1}{x_0}(x-x_0) + \ln(x_0)   \\
			        & = \frac{x}{x_0} - 1 + \ln(x_0)\dots
		      \end{align*}
		\item
	\end{enumerate}
\end{proof}

\end{document}
